%% ma: L11-B16-0071
%% lop: 11
%% chuong: 5
%% bai: 16
%% dang: 11.16.3
%% loai: TLN
%% muc: VD
%% dap_an: "0"
%% nguon: "tai-lieu-hoc-tap-mon-toan-lop-11-theo-chuong-trinh-giao-duc-pho-thong-2018.pdf (Lớp Toán Thầy Duy, Phạm Lê Duy, THPT Chu Văn An), Chương 3 Bài 2, Phần 3, Câu 9 (tr.381)"
%% kien_thuc: "Lớp 11 Bài 16 (giới hạn của hàm số)"
%% trang_thai: chinh-thuc
%% ghi_chu: "Đề gốc ghi 1/x^3 và nhắc 'a,b,c,d' (không có d); kho giữ nguyên chữ của đề, thay mũ 3 bằng x^c để phù hợp 'a,b,c' (đáp số gốc 0 khớp). Đã duyệt 03/11/2026."
\begin{ex}
Cho \[\lim\limits_{x\to+\infty}\dfrac{(4-x^2)^6\cdot(3x^2+2)^2}{(5x^2-3)^2\cdot(1-2x^5)^3}=\lim\limits_{x\to+\infty}\dfrac1{x^c}\cdot\dfrac{\left(\dfrac4{x^2}-1\right)^6\cdot\left(3+\dfrac2{x^2}\right)^b}{\left(5-\dfrac3{x^2}\right)^a\cdot\left(\dfrac1{x^5}-2\right)^3}\] với $a,b,c,d$ là các số tự nhiên. Tính $S=a-b$.
\shortans{0}
\loigiai{\sloppy $(3x^2+2)^2=x^4\left(3+\dfrac2{x^2}\right)^2\Rightarrow b=2$; $(5x^2-3)^2=x^4\left(5-\dfrac3{x^2}\right)^2\Rightarrow a=2$. Vậy $S=a-b=0$ (và $x^{16}:x^{19}=\dfrac1{x^3}$, tức $c=3$).}
\end{ex}
